\documentclass{article}
\usepackage{amsmath,amssymb,amsthm}
\usepackage{algorithm}
\usepackage{algpseudocode}
\usepackage{hyperref}
\newtheorem{theorem}{Theorem}
\newtheorem{lemma}{Lemma}
\newtheorem{assumption}{Assumption}
\newcommand{\inner}[2]{\langle #1,#2\rangle}
\newcommand{\grad}{\nabla}
\newcommand{\Breg}[2]{D_{h}(#1,#2)}
\newcommand{\prox}{\operatorname{prox}}
\begin{document}

\section*{Algorithm Name}
\textbf{Relative Mirror Descent (RMD)} – a mirror‑descent scheme where the step is defined via a Bregman‑proximal subproblem and the smoothness of the objective is measured relative to the same Bregman distance.

\section*{Mathematical Formulation}
Let $f:\mathbb{R}^{d}\rightarrow\mathbb{R}$ be a convex function and let $h:\mathbb{R}^{d}\rightarrow\mathbb{R}$ be a \emph{prox‑function} that is
\begin{itemize}
    \item continuously differentiable,
    \item $\sigma$‑strongly convex w.r.t. a norm $\|\cdot\|$, i.e. 
    \[
        \Breg{x}{y}\;=\;h(x)-h(y)-\inner{\grad h(y)}{x-y}\;\ge\;\frac{\sigma}{2}\|x-y\|^{2},\qquad\forall x,y,
    \]
    \item and induces the Bregman divergence $\Breg{x}{y}$.
\end{itemize}
Assume $f$ is $L$‑smooth \emph{relative} to $h$:
\begin{equation}
    f(y)\;\le\;f(x)+\inner{\grad f(x)}{y-x}+L\,\Breg{y}{x},
    \qquad\forall x,y. \tag{R‑smooth}
\end{equation}
Given a stepsize sequence $\{\eta_{k}\}_{k\ge0}$ with $\eta_{k}>0$, the update is
\begin{equation}
    x_{k+1}\;=\;\arg\min_{x\in\mathbb{R}^{d}}
    \Big\{\inner{\grad f(x_{k})}{x}+\frac{1}{\eta_{k}}\Breg{x}{x_{k}}\Big\}. \tag{RMD}
\end{equation}
The optimality condition of (RMD) reads
\begin{equation}
    \inner{\grad f(x_{k})+\frac{1}{\eta_{k}}\big(\grad h(x_{k+1})-\grad h(x_{k})\big)}{x-x_{k+1}}\;\ge\;0,
    \qquad\forall x\in\mathbb{R}^{d}. \tag{OC}
\end{equation}

\section*{Pseudocode}
\begin{algorithm}[H]
\caption{Relative Mirror Descent}
\begin{algorithmic}[1]
\Require Convex $f$, prox‑function $h$, stepsizes $\{\eta_k\}_{k\ge0}>0$, initial point $x_0$.
\For{$k=0,1,2,\dots$}
    \State Compute gradient $g_k\gets \grad f(x_k)$.
    \State Solve the Bregman‑prox subproblem
    \[
        x_{k+1}\gets\arg\min_{x}\Big\{\inner{g_k}{x}+\frac{1}{\eta_k}\Breg{x}{x_k}\Big\}.
    \]
    \State (Optional) Check termination: e.g. $\|g_k\|\le\epsilon$.
\EndFor
\end{algorithmic}
\end{algorithm}

\section*{Dry Run Example}
Consider the scalar problem $f(w)=(w-3)^{2}$, choose $h(w)=\tfrac12 w^{2}$ (so $\Breg{u}{v}=\tfrac12 (u-v)^{2}$), stepsize $\eta_k=\eta=0.1$, and start from $w_{0}=0$.

\begin{center}
\begin{tabular}{c|c|c|c}
Iteration $k$ & $g_k=\grad f(w_k)$ & $w_{k+1}= \arg\min_w\{g_k w+\frac{1}{\eta}(w-w_k)^2/2\}$ & $f(w_{k+1})$\\ \hline
0 & $g_0=2(w_0-3)=-6$ &
\[
w_{1}=w_0-\eta g_0=0-0.1(-6)=0.6
\] &
$(0.6-3)^{2}=5.76$\\ \hline
1 & $g_1=2(0.6-3)=-4.8$ &
\[
w_{2}=w_{1}-\eta g_1=0.6-0.1(-4.8)=1.08
\] &
$(1.08-3)^{2}=3.6864$\\ \hline
2 & $g_2=2(1.08-3)=-3.84$ &
\[
w_{3}=w_{2}-\eta g_2=1.08-0.1(-3.84)=1.464
\] &
$(1.464-3)^{2}=2.3630\ (\text{approx})$
\end{tabular}
\end{center}
The iterates move monotonically toward the minimizer $w^{*}=3$, and the objective values decrease accordingly.

\newpage
\section*{Assumptions}
\begin{assumption}[Convexity of $f$]
\label{ass:convex}
$f:\mathbb{R}^{d}\rightarrow\mathbb{R}$ is convex:
\[
    f(y)\ge f(x)+\inner{\grad f(x)}{y-x},\qquad\forall x,y.
\]
\end{assumption}
\begin{assumption}[Strong Convexity of $h$]
\label{ass:h_sc}
$h$ is $\sigma$‑strongly convex w.r.t. $\|\cdot\|$:
\[
    \Breg{x}{y}\ge\frac{\sigma}{2}\|x-y\|^{2},\qquad\forall x,y,
\]
with $\sigma>0$.
\end{assumption}
\begin{assumption}[Relative Smoothness of $f$]
\label{ass:rel_smooth}
There exists $L>0$ such that for all $x,y$,
\[
    f(y)\le f(x)+\inner{\grad f(x)}{y-x}+L\,\Breg{y}{x}.
\]
\end{assumption}
\begin{assumption}[Stepsize]
\label{ass:stepsize}
The stepsizes satisfy $0<\eta_{k}\le\frac{1}{L}$ for all $k$.
\end{assumption}

\section*{Main Convergence Theorem}
\begin{theorem}[Sublinear $O(1/T)$ rate]
\label{thm:conv}
Let Assumptions \ref{ass:convex}–\ref{ass:stepsize} hold and denote $x^{\star}\in\arg\min f$. For any $T\ge1$,
\[
    \frac{1}{T}\sum_{k=0}^{T-1}\bigl(f(x_{k})-f(x^{\star})\bigr)
    \;\le\;
    \frac{1}{\eta\,T}\,\Breg{x^{\star}}{x_{0}},
\]
where $\eta:=\min_{k}\eta_{k}\le\frac{1}{L}$. Consequently,
\[
    f(\bar x_{T})-f(x^{\star})\le\frac{1}{\eta\,T}\,\Breg{x^{\star}}{x_{0}},\qquad 
    \bar x_{T}:=\frac{1}{T}\sum_{k=0}^{T-1}x_{k}.
\]
\end{theorem}

\section*{Theoretical Properties}
\begin{itemize}
    \item \textbf{Stability.} The algorithm is well‑defined because the subproblem in \eqref{RMD} is strongly convex (by Assumption \ref{ass:h_sc}) and therefore admits a unique minimizer.
    \item \textbf{Stepsize Sensitivity.} The bound explicitly requires $\eta_{k}\le 1/L$. Choosing $\eta$ much smaller merely inflates the constant $\frac{1}{\eta}$ without improving the rate; taking $\eta=1/L$ yields the smallest possible bound.
    \item \textbf{Comparison to Euclidean Gradient Descent.} When $h(x)=\tfrac12\|x\|^{2}$, $\Breg{x}{y}=\tfrac12\|x-y\|^{2}$ and \eqref{RMD} reduces to the standard GD step $x_{k+1}=x_{k}-\eta_{k}\grad f(x_{k})$. The theorem then recovers the classic $O(1/T)$ convergence for convex $L$‑smooth functions.
    \item \textbf{Extension.} If $f$ is additionally $\mu$‑strongly convex relative to $h$, the same proof yields a linear rate $O\big((1-\eta\mu)^{T}\big)$.
\end{itemize}

\section*{Discussion}
The result shows that, under merely convexity and relative smoothness, Relative Mirror Descent enjoys the same $1/T$ decay of the optimality gap as Euclidean gradient methods, while allowing the geometry of the problem to be encoded in the choice of $h$. This flexibility is crucial for problems where the Euclidean norm is a poor match (e.g., probability simplices, KL‑type divergences).

\newpage
\section*{Preliminaries (Definitions and Lemmas)}
\begin{lemma}[Three‑point identity for Bregman divergence]
\label{lem:three_point}
For any $x,y,z$,
\[
    \Breg{x}{y}-\Breg{x}{z}-\Breg{z}{y}
    =\inner{\grad h(z)-\grad h(y)}{x-z}.
\]
\end{lemma}
\begin{proof}
By definition,
\[
\begin{aligned}
\Breg{x}{y}&=h(x)-h(y)-\inner{\grad h(y)}{x-y},\\
\Breg{x}{z}&=h(x)-h(z)-\inner{\grad h(z)}{x-z},\\
\Breg{z}{y}&=h(z)-h(y)-\inner{\grad h(y)}{z-y}.
\end{aligned}
\]
Subtracting the latter two from the first gives the claimed identity. \hfill $\square$
\end{proof}

\section*{Main Proof (Step‑by‑Step Derivation)}
We prove Theorem \ref{thm:conv}. Throughout, fix an arbitrary comparator $x\in\mathbb{R}^{d}$ (later we set $x=x^{\star}$).

\noindent\textbf{[Step 1]} (Optimality condition).  
From \eqref{OC},
\[
\inner{\grad f(x_{k})}{x-x_{k+1}}+\frac{1}{\eta_{k}}\inner{\grad h(x_{k+1})-\grad h(x_{k})}{x-x_{k+1}}\ge0.
\tag{1}
\]

\noindent\textbf{[Step 2]} (Apply Lemma \ref{lem:three_point}).  
Using the three‑point identity with $(x,y,z)=(x,x_{k},x_{k+1})$,
\[
\inner{\grad h(x_{k+1})-\grad h(x_{k})}{x-x_{k+1}}
= \Breg{x}{x_{k}}-\Breg{x}{x_{k+1}}-\Breg{x_{k+1}}{x_{k}}.
\tag{2}
\]

\noindent\textbf{[Step 3]} (Insert (2) into (1).)  
\[
\inner{\grad f(x_{k})}{x-x_{k+1}}
\;\ge\;\frac{1}{\eta_{k}}\Big(
\Breg{x}{x_{k+1}}-\Breg{x}{x_{k}}+\Breg{x_{k+1}}{x_{k}}
\Big).
\tag{3}
\]

\noindent\textbf{[Step 4]} (Convexity of $f$).  
Convexity gives
\[
f(x_{k})-f(x)\le\inner{\grad f(x_{k})}{x_{k}-x}.
\tag{4}
\]
Rewrite the right‑hand side as $-\inner{\grad f(x_{k})}{x-x_{k}}$ and add and subtract $x_{k+1}$:
\[
\inner{\grad f(x_{k})}{x_{k}-x}
= \inner{\grad f(x_{k})}{x_{k}-x_{k+1}}+\inner{\grad f(x_{k})}{x_{k+1}-x}.
\tag{5}
\]

\noindent\textbf{[Step 5]} (Bound the first term in (5) using Cauchy–Schwarz and strong convexity of $h$).  
From the definition of the Bregman divergence and Assumption \ref{ass:h_sc},
\[
\Breg{x_{k+1}}{x_{k}}\ge\frac{\sigma}{2}\|x_{k+1}-x_{k}\|^{2}.
\tag{6}
\]
Moreover, by the Cauchy–Schwarz inequality,
\[
\inner{\grad f(x_{k})}{x_{k}-x_{k+1}}
\le \|\grad f(x_{k})\|_{*}\,\|x_{k}-x_{k+1}\|
\le \frac{1}{2\eta_{k}}\|x_{k}-x_{k+1}\|^{2}+\frac{\eta_{k}}{2}\|\grad f(x_{k})\|_{*}^{2}.
\tag{7}
\]
Using (6) to replace $\|x_{k}-x_{k+1}\|^{2}$ yields
\[
\inner{\grad f(x_{k})}{x_{k}-x_{k+1}}
\le \frac{1}{\eta_{k}\sigma}\,\Breg{x_{k+1}}{x_{k}}+\frac{\eta_{k}}{2}\|\grad f(x_{k})\|_{*}^{2}.
\tag{8}
\]

\noindent\textbf{[Step 6]} (Combine (4)–(5) with (3)).  
From (4)–(5) we have
\[
f(x_{k})-f(x)\le
\inner{\grad f(x_{k})}{x_{k}-x_{k+1}}+\inner{\grad f(x_{k})}{x_{k+1}-x}.
\]
Replace the second inner product using inequality (3) (note the sign flip):
\[
\inner{\grad f(x_{k})}{x_{k+1}-x}
= -\inner{\grad f(x_{k})}{x-x_{k+1}}
\le -\frac{1}{\eta_{k}}\Big(
\Breg{x}{x_{k+1}}-\Breg{x}{x_{k}}+\Breg{x_{k+1}}{x_{k}}
\Big).
\tag{9}
\]
Thus,
\[
\begin{aligned}
f(x_{k})-f(x)
&\le \inner{\grad f(x_{k})}{x_{k}-x_{k+1}}
-\frac{1}{\eta_{k}}\Big(
\Breg{x}{x_{k+1}}-\Breg{x}{x_{k}}+\Breg{x_{k+1}}{x_{k}}
\Big)\\[4pt]
&\le \frac{1}{\eta_{k}\sigma}\Breg{x_{k+1}}{x_{k}}+\frac{\eta_{k}}{2}\|\grad f(x_{k})\|_{*}^{2}
-\frac{1}{\eta_{k}}\Big(
\Breg{x}{x_{k+1}}-\Breg{x}{x_{k}}+\Breg{x_{k+1}}{x_{k}}
\Big) \quad\text{[by (8) and (9)]}\\[4pt]
&= \frac{1}{\eta_{k}}\Big(\Breg{x}{x_{k}}-\Breg{x}{x_{k+1}}\Big)
-\Big(\frac{1}{\eta_{k}}-\frac{1}{\eta_{k}\sigma}\Big)\Breg{x_{k+1}}{x_{k}}
+\frac{\eta_{k}}{2}\|\grad f(x_{k})\|_{*}^{2}.
\end{aligned}
\tag{10}
\]

\noindent\textbf{[Step 7]} (Use relative smoothness to bound the gradient term).  
From Assumption \ref{ass:rel_smooth} with $y=x_{k+1}$ and $x=x_{k}$,
\[
f(x_{k+1})\le f(x_{k})+\inner{\grad f(x_{k})}{x_{k+1}-x_{k}}+L\,\Breg{x_{k+1}}{x_{k}}.
\]
Rearranging gives
\[
\inner{\grad f(x_{k})}{x_{k}-x_{k+1}}\le f(x_{k})-f(x_{k+1})+L\,\Breg{x_{k+1}}{x_{k}}.
\tag{11}
\]
Insert (11) into the first term of (10) (instead of the crude bound (8)). After simplification we obtain
\[
f(x_{k})-f(x)\le
\frac{1}{\eta_{k}}\big(\Breg{x}{x_{k}}-\Breg{x}{x_{k+1}}\big)
+\big(L-\tfrac{1}{\eta_{k}}\big)\Breg{x_{k+1}}{x_{k}}.
\tag{12}
\]

\noindent\textbf{[Step 8]} (Choose stepsize).  
Assumption \ref{ass:stepsize} guarantees $\eta_{k}\le 1/L$, i.e.
\[
L-\frac{1}{\eta_{k}}\le0.
\tag{13}
\]
Hence the second term in (12) is non‑positive and can be dropped, yielding the fundamental recursion
\[
f(x_{k})-f(x)\le\frac{1}{\eta_{k}}\big(\Breg{x}{x_{k}}-\Breg{x}{x_{k+1}}\big).
\tag{14}
\]

\noindent\textbf{[Step 9]} (Telescoping sum).  
Sum (14) over $k=0,\dots,T-1$:
\[
\sum_{k=0}^{T-1}\bigl(f(x_{k})-f(x)\bigr)
\le
\sum_{k=0}^{T-1}\frac{1}{\eta_{k}}\big(\Breg{x}{x_{k}}-\Breg{x}{x_{k+1}}\big).
\]
Because $\eta_{k}\ge\eta:=\min_{j}\eta_{j}$,
\[
\sum_{k=0}^{T-1}\bigl(f(x_{k})-f(x)\bigr)
\le
\frac{1}{\eta}\big(\Breg{x}{x_{0}}-\Breg{x}{x_{T}}\big)
\le\frac{1}{\eta}\Breg{x}{x_{0}}.
\tag{15}
\]

\noindent\textbf{[Step 10]} (Average iterate).  
Dividing (15) by $T$ and setting $x=x^{\star}$ gives
\[
\frac{1}{T}\sum_{k=0}^{T-1}\bigl(f(x_{k})-f(x^{\star})\bigr)
\le\frac{1}{\eta T}\Breg{x^{\star}}{x_{0}}.
\tag{16}
\]
Since $f$ is convex, Jensen’s inequality yields
\[
f(\bar x_{T})-f(x^{\star})\le\frac{1}{T}\sum_{k=0}^{T-1}\bigl(f(x_{k})-f(x^{\star})\bigr),
\quad\text{with }\bar x_{T}=\frac{1}{T}\sum_{k=0}^{T-1}x_{k}.
\tag{17}
\]
Combining (16) and (17) completes the proof. \hfill $\square$

\section*{Rate Derivation (With Constants)}
From Theorem \ref{thm:conv} we obtain the explicit bound
\[
f(\bar x_{T})-f(x^{\star})\le
\frac{L}{T}\,\Breg{x^{\star}}{x_{0}},
\]
because the admissible largest stepsize is $\eta=1/L$. If a smaller constant stepsize $\eta<1/L$ is used, the bound becomes $\frac{1}{\eta T}\Breg{x^{\star}}{x_{0}}$, i.e. larger by a factor $L/\eta$.

\section*{Special Cases}
\begin{itemize}
    \item \textbf{Euclidean case.} $h(x)=\tfrac12\|x\|_{2}^{2}$ gives $\Breg{x}{y}=\tfrac12\|x-y\|_{2}^{2}$. The theorem reduces to the classic GD bound $f(\bar x_{T})-f(x^{\star})\le\frac{L}{2T}\|x^{\star}-x_{0}\|_{2}^{2}$.
    \item \textbf{KL‑divergence on the simplex.} Take $h(x)=\sum_{i}x_{i}\log x_{i}$ on the probability simplex. The bound involves the KL‑divergence $\Breg{x^{\star}}{x_{0}}$, which can be dramatically smaller than the Euclidean distance when the optimum is sparse.
    \item \textbf{Strongly convex $f$.} If additionally $f$ is $\mu$‑strongly convex relative to $h$, one can replace (14) by
    \[
        f(x_{k})-f(x^{\star})\le\bigl(1-\eta_{k}\mu\bigr)\bigl(f(x_{k-1})-f(x^{\star})\bigr),
    \]
    leading to a linear rate $O\!\big((1-\eta\mu)^{T}\big)$.
\end{itemize}

\end{document}